\makeatletter
\def\input@path{{./}{../}{../../}{../../../}{../../../../}{preamble/}{../preamble/}{../../preamble/}{../../../preamble/}}
\makeatother
\input{preamble.tex}

\title{Grade 11 -- Term 1 \\ C8 Practice Activity with Solutions}
\author{}
\date{}

\begin{document}
\maketitle
\setcounter{tocdepth}{1}
\setcounter{secnumdepth}{2}
\phantomsection
\label{toc}
\tableofcontents

\section*{Evaluated criterion}
\begin{itemize}
\item[\textbf{C8.}] Uses Student's $t$-distribution tables, degrees of freedom, and one-tailed or two-tailed areas to determine critical values, central probabilities, and tail probabilities or probability ranges.
\end{itemize}
\vspace{-4mm}

\section{Introduction to Student's \texorpdfstring{$t$}{t}-Distribution}
\label{sec:c8-introduction}
\SectionProblemsTableOfContents{
\SectionProblemLink{sec:c8-model}{Shape, notation, and degrees of freedom} \\
\SectionProblemLink{sec:c8-areas}{Tail areas, symmetry, and critical values} \\
\SectionProblemLink{sec:c8-table}{Reading the Student's \texorpdfstring{$t$}{t}-table} \\
\SectionProblemLink{sec:c8-examples}{Worked table examples}
}
\vspace{5mm}

\subsection{Shape, notation, and degrees of freedom}
\label{sec:c8-model}
Student's $t$-distribution is a family of continuous probability distributions used to model standardized sample means when the population standard deviation is unknown and is replaced by the sample standard deviation. Write
\[
T\sim t_\nu,
\]
where $\nu$ (the Greek letter nu) is the number of \textbf{degrees of freedom}. Each curve is bell-shaped and symmetric about zero, with total area $1$. Compared with the standard normal distribution, a $t$-distribution has heavier tails: large positive and negative standardized values are more likely.

For an independent random sample of $n$ observations from a normal population, the standardized sample mean
\[
T=\frac{\bar X-\mu}{S/\sqrt n}\sim t_{n-1}.
\]
Here $\bar X$ is the sample mean, $\mu$ the population mean, and $S$ the sample standard deviation. Estimating the mean leaves $n-1$ independent deviations, so $\nu=n-1$. This rule applies to the sample-mean models in this activity; degrees of freedom need not equal $n-1$ in every statistical model. No calculation of $T$ from raw data is required here.

Small degrees of freedom give heavier tails and larger positive critical values for the same upper-tail area. As $\nu$ increases, the curve approaches the standard normal curve. A $t$-value is dimensionless; it is not an expenditure, return, or delivery time in the original units.
\secInicio[sec:c8-model][sec:c8-introduction]

\subsection{Tail areas, symmetry, and critical values}
\label{sec:c8-areas}
For $t>0$, let $q=P(T>t)$. Symmetry and complementary areas give
\[
\begin{aligned}
P(T<-t)&=P(T>t)=q,\\
P(T<t)&=1-q,\\
P(|T|>t)&=P(T<-t)+P(T>t)=2q,\\
P(-t<T<t)&=1-2q.
\end{aligned}
\]
Thus a \textbf{one-tailed} probability refers to one end of the distribution. A \textbf{two-tailed} probability combines both ends outside $-t$ and $t$. A \textbf{central} probability is the area between these two cutoffs. For any real $a$, symmetry gives $P(T<a)=P(T>-a)$. Since the distribution is continuous, including or excluding an endpoint does not change its probability.

A \textbf{critical $t$-value} is a cutoff that leaves a specified area in a tail or between symmetric cutoffs. In this document, $t_{\nu,q}>0$ denotes the cutoff satisfying $P(T>t_{\nu,q})=q$:
\[
\begin{array}{c|c}
\text{Required area} & \text{Cutoff or cutoffs}\\ \hline
P(T>t)=\alpha & t=t_{\nu,\alpha}\\
P(T<c)=\alpha\quad (\alpha<0.5) & c=-t_{\nu,\alpha}\\
P(|T|>t)=\alpha & -t_{\nu,\alpha/2},\ t_{\nu,\alpha/2}\\
P(-t<T<t)=p & -t_{\nu,(1-p)/2},\ t_{\nu,(1-p)/2}
\end{array}
\]
\secInicio[sec:c8-areas][sec:c8-introduction]

\subsection{Reading the Student's \texorpdfstring{$t$}{t}-table}
\label{sec:c8-table}
Each entry below is a positive $t_{\nu,q}$; the column heading $q$ is the \textbf{upper-tail probability}, not the central area or the combined two-tail area. First choose the row for $\nu$, then convert the requested area into $q$, and finally read the entry. The extra header rows show the corresponding two-tail and central areas.

\begin{center}
\begingroup
\setlength{\tabcolsep}{9pt}
\begin{tabular}{c|rrrrr}
\toprule
Upper tail $q$ & $0.10$ & $0.05$ & $0.025$ & $0.01$ & $0.005$\\
Two tails $2q$ & $0.20$ & $0.10$ & $0.05$ & $0.02$ & $0.01$\\
Central $1-2q$ & $0.80$ & $0.90$ & $0.95$ & $0.98$ & $0.99$\\
\midrule
Degrees of freedom $\nu$ & \multicolumn{5}{c}{Positive critical value $t_{\nu,q}$}\\
\midrule
5  & 1.476 & 2.015 & 2.571 & 3.365 & 4.032\\
9  & 1.383 & 1.833 & 2.262 & 2.821 & 3.250\\
11 & 1.363 & 1.796 & 2.201 & 2.718 & 3.106\\
14 & 1.345 & 1.761 & 2.145 & 2.624 & 2.977\\
19 & 1.328 & 1.729 & 2.093 & 2.539 & 2.861\\
24 & 1.318 & 1.711 & 2.064 & 2.492 & 2.797\\
29 & 1.311 & 1.699 & 2.045 & 2.462 & 2.756\\
\bottomrule
\end{tabular}
\endgroup
\end{center}

Entries are rounded to three decimal places. Report critical values to this precision. When a question uses a printed entry as its cutoff, report the associated probability as a table-based approximation. If a cutoff lies strictly between two entries, give a probability interval: larger positive cutoffs leave smaller upper-tail areas. Do not interpolate or claim an exact probability between columns. Every sample size used below has its required row in this table.
\secInicio[sec:c8-table][sec:c8-introduction]

\subsection{Worked table examples}
\label{sec:c8-examples}
\subsubsection*{C8 Example 1 -- A positive and a negative cutoff}
An analyst models a standardized mean return from $n=10$ independent observations using $T\sim t_{n-1}$. Find $a$ with $P(T>a)=0.05$ and $b$ with $P(T<b)=0.05$.

The degrees of freedom are $\nu=10-1=9$. Each request concerns one tail. Row $9$, column $q=0.05$, gives $1.833$. Therefore,
\[
a\approx1.833,\qquad b\approx-1.833.
\]
The negative cutoff follows from $P(T<-1.833)=P(T>1.833)\approx0.05$.

\subsubsection*{C8 Example 2 -- Two tails and a central area}
A standardized household-expenditure mean uses $n=15$ observations, so $T\sim t_{14}$. Find symmetric cutoffs enclosing $0.95$ of the distribution.

The area outside is $1-0.95=0.05$, shared equally between the two tails. Use $q=0.05/2=0.025$, not $0.05$. Row $14$ gives
\[
t_{14,0.025}\approx2.145,\qquad
P(-2.145<T<2.145)\approx0.95.
\]
The cutoffs are approximately $-2.145$ and $2.145$; equivalently, $P(|T|>2.145)\approx0.05$.

\subsubsection*{C8 Example 3 -- A probability range and symmetry}
A standardized delivery-time mean has $n=20$, so $T\sim t_{19}$. Bound $P(T>2.30)$, $P(T<-2.30)$, and $P(|T|>2.30)$.

In row $19$, $2.093<2.30<2.539$. The bounding entries belong to $q=0.025$ and $q=0.01$, respectively. Thus
\[
0.01<P(T>2.30)<0.025,\qquad
0.01<P(T<-2.30)<0.025.
\]
Doubling the equal tail areas and then taking the complement gives
\[
0.02<P(|T|>2.30)<0.05,\qquad
0.95<P(-2.30<T<2.30)<0.98.
\]
\secInicio[sec:c8-examples][sec:c8-introduction]

\section{Calculating Probabilities and Critical Values Using a \texorpdfstring{$t$}{t}-Table}
\label{sec:c8-practice}
For each problem, record the degrees of freedom, classify the requested area, and identify the upper-tail column or the two columns bounding it. All sample-mean models below assume independent random observations from a normal population, with unknown population standard deviation replaced by the sample standard deviation; therefore $T\sim t_{n-1}$. Use only the included table. Solutions follow each problem for checking after an independent attempt.
\SectionProblemsTableOfContents{
\SectionProblemLink{sec:c8-returns}{Investment returns -- right-tail cutoff} \\
\SectionProblemLink{sec:c8-expenditure}{Household expenditures -- left-tail cutoff} \\
\SectionProblemLink{sec:c8-delivery}{Delivery times -- two-tailed cutoffs} \\
\SectionProblemLink{sec:c8-costs}{Production costs -- central cutoffs} \\
\SectionProblemLink{sec:c8-wages}{Wages -- reading a printed entry} \\
\SectionProblemLink{sec:c8-sales}{Sales -- bounding a right-tail probability} \\
\SectionProblemLink{sec:c8-revenues}{Business revenues -- bounding two tails} \\
\SectionProblemLink{sec:c8-prices}{Prices -- interpreting a negative value} \\
\SectionProblemLink{sec:c8-small}{Small business costs -- reading verbal areas} \\
\SectionProblemLink{sec:c8-payroll}{Payroll records -- determining the sample size} \\
\SectionProblemLink{sec:c8-planning}{Revenue planning -- mixed probability ranges}
}
\vspace{5mm}
\subsection{Investment returns -- right-tail cutoff}
\label{sec:c8-returns}
A financial analyst uses $n=10$ daily investment returns. The standardized sample mean has $T\sim t_9$, with $\nu=9$.

\QuestionsTasks[sec:c8-practice]{{8}/{sol:c8-returns}/{Find the positive cutoff $t$ such that $P(T>t)=0.10$.}}

\subsubsection*{C8 Solution}
\phantomsection
\label{sol:c8-returns}
This is one right tail. Use row $9$ and column $q=0.10$: $t\approx1.383$.

\secInicio[sec:c8-returns][sec:c8-practice]

\subsection{Household expenditures -- left-tail cutoff}
\label{sec:c8-expenditure}
A survey uses $n=12$ household expenditures. The standardized mean has $T\sim t_{11}$, with $\nu=11$.

\QuestionsTasks[sec:c8-practice]{{8}/{sol:c8-expenditure}/{Find the negative cutoff $c$ such that $P(T<c)=0.025$.}}

\subsubsection*{C8 Solution}
\phantomsection
\label{sol:c8-expenditure}
By symmetry, use the upper-tail column $q=0.025$ in row $11$, then change the sign: $c\approx-2.201$.

\secInicio[sec:c8-expenditure][sec:c8-practice]

\subsection{Delivery times -- two-tailed cutoffs}
\label{sec:c8-delivery}
A logistics company uses delivery times from $n=15$ deliveries to model a standardized sample mean.

\QuestionsTasks[sec:c8-practice]{{8}/{sol:c8-delivery}/{Determine $\nu$ and find $t>0$ such that $P(|T|>t)=0.10$. State both cutoffs and the area in each tail.}}

\subsubsection*{C8 Solution}
\phantomsection
\label{sol:c8-delivery}
$\nu=15-1=14$. The combined two-tail area is $0.10$, so each tail has $q=0.05$. Row $14$ gives $t\approx1.761$; the cutoffs are $-1.761$ and $1.761$.

\secInicio[sec:c8-delivery][sec:c8-practice]

\subsection{Production costs -- central cutoffs}
\label{sec:c8-costs}
A manufacturer uses production costs from $n=20$ independent batches to model a standardized mean.

\QuestionsTasks[sec:c8-practice]{{8}/{sol:c8-costs}/{Find symmetric cutoffs $-t$ and $t$ enclosing a central probability of $0.98$. State the area outside them.}}

\subsubsection*{C8 Solution}
\phantomsection
\label{sol:c8-costs}
$\nu=19$. The outside area is $0.02$, so $q=0.01$. Row $19$ gives cutoffs $-2.539$ and $2.539$.

\secInicio[sec:c8-costs][sec:c8-practice]

\subsection{Wages -- reading a printed entry}
\label{sec:c8-wages}
A business models a standardized mean wage using a sample of $n=25$ employees.

\QuestionsTasks[sec:c8-practice]{{8}/{sol:c8-wages}/{Use the table to estimate $P(T>2.064)$ and $P(T<2.064)$.}}

\subsubsection*{C8 Solution}
\phantomsection
\label{sol:c8-wages}
$\nu=24$. The entry $2.064$ is in column $q=0.025$. Therefore $P(T>2.064)\approx0.025$ and $P(T<2.064)\approx0.975$.

\secInicio[sec:c8-wages][sec:c8-practice]

\subsection{Sales -- bounding a right-tail probability}
\label{sec:c8-sales}
A retailer uses sales from $n=30$ randomly selected trading days to model a standardized mean.

\QuestionsTasks[sec:c8-practice]{{8}/{sol:c8-sales}/{Give a probability interval for $P(T>1.90)$, identifying the two table entries used.}}

\subsubsection*{C8 Solution}
\phantomsection
\label{sol:c8-sales}
$\nu=29$. Since $1.699<1.90<2.045$, the relevant columns are $0.05$ and $0.025$. Hence $0.025<P(T>1.90)<0.05$.

\secInicio[sec:c8-sales][sec:c8-practice]

\subsection{Business revenues -- bounding two tails}
\label{sec:c8-revenues}
A revenue model uses standardized sample means from $n=12$ independent businesses.

\QuestionsTasks[sec:c8-practice]{{8}/{sol:c8-revenues}/{Give probability intervals for $P(|T|>2.50)$ and $P(-2.50<T<2.50)$.}}

\subsubsection*{C8 Solution}
\phantomsection
\label{sol:c8-revenues}
$\nu=11$. Since $2.201<2.50<2.718$, $0.01<P(T>2.50)<0.025$. Doubling gives $0.02<P(|T|>2.50)<0.05$. The central area therefore satisfies $0.95<P(-2.50<T<2.50)<0.98$.

\secInicio[sec:c8-revenues][sec:c8-practice]

\subsection{Prices -- interpreting a negative value}
\label{sec:c8-prices}
A price study models standardized mean prices from samples of $n=15$ shops.

\QuestionsTasks[sec:c8-practice]{{8}/{sol:c8-prices}/{Bound $P(T<-2.30)$ and $P(T>-2.30)$. Explain the role of symmetry.}}

\subsubsection*{C8 Solution}
\phantomsection
\label{sol:c8-prices}
$\nu=14$. The entries $2.145<2.30<2.624$ imply $0.01<P(T>2.30)<0.025$. Symmetry gives $0.01<P(T<-2.30)<0.025$. Taking the complement gives $0.975<P(T>-2.30)<0.99$.

\secInicio[sec:c8-prices][sec:c8-practice]

\subsection{Small business costs -- reading verbal areas}
\label{sec:c8-small}
A standardized mean operating-cost model uses $n=6$ businesses.

\QuestionsTasks[sec:c8-practice]{{8}/{sol:c8-small}/{Find (a) the positive cutoff exceeded in $5\%$ of model outcomes; (b) the symmetric cutoffs with $5\%$ of outcomes outside them in total; and (c) the symmetric cutoffs enclosing the middle $80\%$.}}

\subsubsection*{C8 Solution}
\phantomsection
\label{sol:c8-small}
$\nu=5$. (a) One right tail: $q=0.05$, so $t\approx2.015$. (b) Two tails: $q=0.05/2=0.025$, so the cutoffs are $\pm2.571$. (c) Central area: $q=(1-0.80)/2=0.10$, so the cutoffs are $\pm1.476$.

\secInicio[sec:c8-small][sec:c8-practice]

\subsection{Payroll records -- determining the sample size}
\label{sec:c8-payroll}
A wage study selects $5$ employees from each of $5$ departments. Treat all $25$ observations as independent; the standardized mean follows the sample-mean model described above.

\QuestionsTasks[sec:c8-practice]{{8}/{sol:c8-payroll}/{Find the negative cutoff below which $1\%$ of model outcomes lie. Then find the symmetric cutoffs enclosing the middle $90\%$.}}

\subsubsection*{C8 Solution}
\phantomsection
\label{sol:c8-payroll}
$n=5\times5=25$, so $\nu=24$. For the left tail, use $q=0.01$ and symmetry: $c\approx-2.492$. For the central area, $q=(1-0.90)/2=0.05$, giving cutoffs $\pm1.711$.

\secInicio[sec:c8-payroll][sec:c8-practice]

\subsection{Revenue planning -- mixed probability ranges}
\label{sec:c8-planning}
A standardized mean-revenue model uses samples of $n=10$ stores.

\QuestionsTasks[sec:c8-practice]{{8}/{sol:c8-planning}/{For a cutoff magnitude of $3.00$, bound (a) the probability below $-3.00$; (b) the probability outside the symmetric cutoffs; and (c) the probability between them. Also estimate $P(-1.833<T<2.262)$ from the table.}}

\subsubsection*{C8 Solution}
\phantomsection
\label{sol:c8-planning}
$\nu=9$. Since $2.821<3.00<3.250$, $0.005<P(T>3.00)<0.01$. Thus (a) $0.005<P(T<-3.00)<0.01$; (b) $0.01<P(|T|>3.00)<0.02$; and (c) $0.98<P(-3.00<T<3.00)<0.99$. For the asymmetric interval, the left excluded tail is approximately $0.05$ and the right excluded tail is approximately $0.025$. Therefore $P(-1.833<T<2.262)\approx1-0.05-0.025=0.925$.

\secInicio[sec:c8-planning][sec:c8-practice]

\end{document}


